\chapter{Individual variability in prosody production and perception} \label{ch:background variability}

Inter-individual variability has only recently become a popular topic in prosodic research, although it has always been present in the data. In a relatively early observation, \citet{beckmanArticulatoryEvidenceDifferentiating1994} noted that nuclear accented syllables are longer than unaccented syllables for some but not all speakers of American English. Prosody, much like the segmental layer of speech, is thus characterized by a lack of invariance, missing a consistent one-to-one correspondence between features of the speech signal and phonological categories \citep{perkellInvarianceVariabilitySpeech1986, niebuhrPerzeptionUndKognitive2007, kurumadaThinkingProbabilisticallyStudy2021,  arvanitiVariabilityOverlapCue2024}. Lack of invariance is also evident in the probabilistic association of prosody and meaning \citep[e.g.,][]{hirschbergCommunicationProsodyFunctional2002, calhounCentralityMetricalStructure2010}. AM theory formalizes this potential for variability in the proposition that metrical structure need not be realized in any one specific phonetic way \citep[55]{laddIntonationalPhonology2008}. Much of the variance in the prosodic signal may in fact be due to individual differences \citep{arvanitiVariabilityOverlapCue2024}. 

This chapter deals with the question of how and why language users differ in their production and perception of prosody. As such, the main concern is \textit{inter}-individual variability. However, individuals may also differ within their own production and perception. Such \textit{intra}-individual variability is often related to register or style \citep[for a review and systematization of intra-individual variability, see][]{ulbrichWhatIntraindividualVariation2021}. Intra-individual variability in the sense of flexibility to adapt is also what allows listeners to deal with speaker-specific variability. Nevertheless, there is evidence that individuals frequently differ systematically from each other but are remarkably consistent within their own production and perception \citep{ouyangIndividualDifferencesProsodic2015, giovannoneIndividualDifferencesLexical2023}.

As established in Section \ref{section: prominence form}, prosodic prominence is redundantly encoded in speech via changes in fundamental frequency (F0), duration, and intensity. Since there are many potential prosodic cues to prominence, it lends itself especially readily to substantial variability. It seems inevitable that speakers may choose different prosodic parameters to encode prominence and that listeners attend to different cues when decoding prominence. In addition, prosody is multi-functional, encoding different aspects of meaning, which could elicit further individual variability depending on the communicative intentions of the speaker \citep{coleProsodyContextReview2015}. Many studies note the presence of such inter-individual differences \citep{baumannTonalArticulatoryMarking2007, griceTonalAssociationDerived2009, muckeEffectFocusMarking2014, griceIntegratingDiscretenessContinuity2017, roessigDynamicsIntonationCategorical2019}, yet few make it their primary object of research \citep{moProsodyProductionPerception2010, cangemiListenerSpecificPerceptionSpeakerSpecific2015, ouyangIndividualDifferencesProsodic2015}. The following sections will summarize pertinent findings from both incidental and systematic observations regarding production (Section \ref{sec:ind_production}) and perception (Section \ref{sec:ind_perception}) before turning to discuss potential reasons for individual variability (Section \ref{sec:causes_variability}).

\section{The nature of individual variability}

\subsection{Production} \label{sec:ind_production}

Previous studies have observed considerable individual variability in production in many areas of prosodic research. Individual differences emerge, for example, in how and to what extent prosodic categories such as pitch accent types, lexical tones, and word stress are distinguished and how and to what extent pragmatic contrasts, e.g., focus and information status, are encoded distinctly via prosody.

The boundaries of prosodic categories are often described as ``fuzzy" and that this is the case is particularly evident in the challenge of prosodic annotation \citep[e.g.,][]{coleNewMethodsProsodic2016}. One of the potential reasons underlying the fuzziness of prosodic categories is inter-speaker variability \citep[e.g.,][]{zahner-ritterThreeKindsRisingFalling2022}. For example, \citet{niebuhrAreThereShapers2011} found that Standard Northern German speakers varied in which prosodic parameters they used to distinguish H* from H+L* accents. Some speakers relied more on peak alignment, referred to as ``aligners", others on the shape of the pitch contour (measured as rise duration divided by fall duration), so-called ``shapers".  Similarly, \citet{grylliaManyShapes2022} observed that Greek speakers varied considerably in how they produced H* accents in terms of F0 contour shape and scaling. They note that this variability could potentially lead to category overlap with other accent types in the Greek inventory, such as L+H* and H*+L.

In extreme cases, speakers may not even produce certain category distinctions because they are not part of their inventories. In an imitation task using stimuli where the F0 peak was incrementally adjusted from a medial to a late peak in 15 steps, \citet{pierrehumbertCategoriesTonalAlignment1989} found that four out of five American English speakers produced two types of pitch accents, consistent with L+H* (medial peak) and L*+H (late peak), although the location of the category boundary differed across speakers. However, they also observed that one speaker did not produce any instances of late peaks, which they interpreted as either a strong personal preference for medial peaks or a lack of an L*+H accent in that speaker's inventory. 

Similar amounts of variability are present in the production of word stress and lexical tones. In a study on word stress production in Dutch, \citet{severijnenYourVOORnaamNot2024} observed that speakers differed in which cues (F0, duration or intensity) they prioritized to encode stress. They also found evidence for cue-trading, in that speakers who predominantly relied on F0 to encode stress typically relied less on intensity. However, they noted a positive correlation between intensity and duration: If a speaker used intensity as a stress correlate, they also likely made use of duration. This finding lends further validity to measures integrating both dimensions \citep{beckmanStressNonStressAccent1986, albertModelSonorityBased2023}. There is also evidence for individual variability in the production of lexical tones, for example in Cantonese speakers \citep{wongPerceptualNormalizationInter2003, chanSpeakerVariabilityRealisation2016}. \citet{chanSpeakerVariabilityRealisation2016} observed that speakers varied considerably in the F0 range they produced, alignment of turning points and how well they distinguished different lexical tones, with some speakers even merging the two rising tones. Such a merger may be indicative of an ongoing sound change, which has previously been reported for Cantonese tones \citep{mokProductionPerceptionSound2013}.

Comparable to the fuzziness of prosodic categories, prosody-meaning mapping has often been described as probabilistic \citep[e.g.,][]{calhounCentralityMetricalStructure2010}. An area dealing with the interface of prosody and meaning that has received much attention is the prosodic marking of focus. Although such studies rarely explicitly set out to investigate individual variability, many have incidentally noted that individual speakers may differ in how they prosodically distinguish between different focus types, such as broad, narrow, and contrastive focus. For example, \citet{muckeEffectFocusMarking2014} observed speaker-specific differences in the use of pitch accent (types) and duration as cues to distinguish focus from non-focus and between focus types in five speakers of German. More variability emerged in the latter comparisons: Target words in the background condition were always deaccented and typically shorter than focused target words across speakers, while not all speakers used different pitch accent types and duration to encode the distinction between broad, narrow, and contrastive focus. In fact, one out of the five speakers almost exclusively produced H* accents across all three focus conditions. Taking a closer look at the same dataset, \citet{cangemiListenerSpecificPerceptionSpeakerSpecific2015} found that German speakers differed in the type and number of prosodic cues they employed to encode the three-way contrast between the different focus types. While four out of five speakers used multiple cues, e.g., peak alignment, peak height, and word duration, to produce robustly distinct exponents of each focus type, one speaker used only single cues to differentiate between two focus types, e.g., word duration to distinguish broad from narrow and contrastive focus and the presence or absence of prenuclear accents to distinguish broad and narrow from contrastive focus. 

Similarly, \citet{baumannTonalArticulatoryMarking2007} observed that German speakers differed in which prosodic parameters they used to distinguish broad and contrastive focus, some speakers using more cues than others. For example, five out of six speakers produced later peaks in contrastive than in broad focus, three out of these five speakers additionally produced higher F0 peaks in contrastive focus but one speaker used neither of these cues and in fact distinguished between the two focus types only minimally. Very much in line with these findings, \citet{kuglerProductionPerceptionContrast2015} investigated the production of broad and contrastive focus in eight German speakers and found that only some speakers used higher F0 scaling as a cue to contrastive focus. However, all speakers produced later peak alignment and longer duration in contrastive than in broad focus.

\citet{griceIntegratingDiscretenessContinuity2017} looking at continuous operationalizations of the F0 contour in the same data as \citet{muckeEffectFocusMarking2014} and \citet{cangemiListenerSpecificPerceptionSpeakerSpecific2015} found that while all speakers exhibited the same tendencies in the mapping of peak alignment to focus type, i.e., broad focus was produced with earlier peaks, narrow and contrastive focus with later peaks, speakers differed in the extent to which they distinguished these categories. The same was true for tonal onglide, with broad focus exhibiting the smallest and contrastive focus the largest onglide values across speakers, but individuals differing in how distinct the association between onglide and focus types were. In line with these findings using a different dataset, \citet{roessigDynamicsIntonationCategorical2019} were able to identify two groups of speakers with respect to how strongly they encoded focus types via tonal onglide: The slightly larger group (16 speakers) made smaller differences between the focus types producing predominantly positive onglides across all focus conditions but still the smallest values in broad focus and the largest ones in contrastive focus. The smaller group (11 speakers) made larger distinctions between the focus types, which were also of a categorical nature: They produced mostly negative onglides in broad focus and mostly positive onglides in contrastive focus, which would most likely map onto different pitch accent types.

Speakers may also vary in their preferences for certain pitch accent types in specific meaning conditions. \citet{griceTonalAssociationDerived2009} found that German speakers differed in the extent to which they favored early peaks (H+!H* or H+L*) over plain downstep (!H*) to mark accessible referents: While all speakers preferably produced early peaks, five out of six speakers also produced downsteps, however varying considerably with regard to the proportion of downsteps in their speech. One speaker exclusively produced early peaks and no downsteps.

As briefly noted in Section \ref{sec: focus background}, it is not only the focus exponent that is prosodically marked to indicate focus types. Preceding stretches of speech may also take on a distinct prosodic form, which adds another site for potential speaker differences. \citet{ipSearchSalienceFocus2021} found that speakers of Australian English differed in how many and which cues they used on a preceding word to mark focus. In their sample of four speakers, one speaker modified duration, F0, and intensity cues, another only mean F0 and intensity, a third only pre-target interval duration, and the last duration and intensity.

Not only West-Germanic languages are rife with speaker-specific variability in focus marking. \citet{heldnerTemporalEffectsFocus2001} looked at duration as a cue to focus in six speakers of Swedish. While all speakers lengthened words in focus to some degree, they differed considerably with regard to the extent of this lengthening. The two speakers at the extreme ends of the scale lengthened focal relative to non-focal words by 5\% and 45\%, respectively.
Similarly, in a study including four speakers of Arabic, \citet{dejongComparingStressLexical2002} found that the encoding of focus differed considerably across speakers, while the encoding of stress was more similar. They attribute this difference to a higher degree of conventionalization in stress than in focus in the language.

\citet{cangemiSpeakerspecificIntonationalMarking2016} observed that speakers of Egyptian Arabic differed in whether and to what extent they used F0 turning point alignment and scaling to encode broad versus narrow focus. Some speakers produced differences predominantly in terms of scaling, others in alignment, and yet again others varied both parameters across the two focus types. However, using the same dataset, \citet{bruggemanUnifyingSpeakerVariability2017} found that TCoG (alignment) provided a more uniform characterization of individual speaker behavior: All speakers produced earlier TCoGs in narrow than in broad focus. TCoG likely functions as a reconciling parameter in this data because it is a more holistic measure that weighs the F0 contour over time and thus takes into account both scaling and overall shape of the F0 contour.

Still dealing with focus but widening the scope slightly, \citet{ouyangIndividualDifferencesProsodic2015} investigated the encoding of informativity (a combination of focus type, contextual probability, and word frequency) via the F0 contour in American English. Overall, about half of the 16 speakers in their sample appeared to follow a general pattern of how informativity modulated F0 range, i.e., F0 range was smaller in broad than in narrow focus for frequent and predictable words, while the other half did not produce any systematic alternative patterns. In addition, they observed that some speakers made larger differences in terms of F0 excursion between the levels of the three informativity factors. In a study including three speakers of Dutch, \citet{swertsProsodyMarkerInformation1994} investigated prosodic cues to discourse structuring in unscripted monologues. They found that speakers differed in how consistently they prosodically marked the end of informational units and newly introduced information via pitch movement and intonational boundaries. 

Several studies thus attest to individual differences in focus prosody in various languages. Other aspects of prosody-meaning mapping have received less attention, specifically with regard to individual variability.
\citet{moProsodyProductionPerception2010} investigated correlates of prominence in a corpus of conversational American English speech produced by 25 speakers. Prominence had been annotated in a Rapid Prosody Transcription task by 11 to 20 transcribers per speech excerpt. \citet{moProsodyProductionPerception2010} found that speakers varied in which prosodic parameters they used to encode prominence and how effectively they used these cues. Some speakers relied mostly on intensity, others on F0. Duration was only used by a small number of speakers. Similarly, \citet{dahanInterspeakerVariabilityEmphatic1996} in a study with four speakers of French found inter-individual differences in the production of emphasis. Emphasis was elicited by asking the speakers to ``insist" on a specific target word \citet[345]{dahanInterspeakerVariabilityEmphatic1996}. Speakers varied mostly in the extent to which they increased duration, intensity, and F0 range on the target word and attenuated the same parameters in the preceding and following syllables.

Apart from prominence production, other prosodic domains may also be subject to speaker differences. In a large study with 90 speakers of Southern British English, \citet{peppeProsodicVariationSouthern2000} observed individual differences in the use of lengthening, pausing, pitch reset, and pitch movement to chunk utterances. Considerable variability with regard to the prosodic cues to focus was also observed in the same study. \citet{xieEncodingDecodingMeaning2021} investigated the prosodic marking of questions via final F0 movement and duration in 65 speakers of American English. They found that most speakers employed final F0 rises but some also used final syllable duration. Finally, lots of variability in non-canonical word order marking in German was observed by \citet{huttenlauchIndividualVariabilityProsodic2022}. Only two out of 16 speakers produced longer pauses after the noun in OVS sentences than in SVO sentences. In addition, ten speakers differentiated between SVO and OVS sentences via F0, mostly with a larger range in SVO than in OVS sentences.

In contrast to all the findings on variability, \citet{yoonSpeakerConsistencyRealization2010} emphasizes similarities across speakers: In a corpus of read speech, in which five radio announcers (speakers of American English) produced the same texts, \citet{yoonSpeakerConsistencyRealization2010} found agreement rates of 80\% on the presence of a pitch accent and 72\% for type of pitch accent produced. These speakers were of course highly trained in a very specific speech style, which could contribute to the apparent uniformity.

To summarize, there is ample evidence for inter-speaker variability in prosody production in many languages, pertaining to the sharpness of prosodic category distinctions and robustness of prosody-meaning mapping, as well as to the type of prosodic parameters that are used to encode both. However, most observations are based on very small speaker samples and/or studies on focus. Nevertheless, the pervasive presence of inter-speaker variability in the encoding of several functional contrasts poses challenges for the listener. The next section will thus address individual variability from the perception side.

\subsection{Perception} \label{sec:ind_perception}

Two questions need to be addressed from a perception perspective: First, how do listeners deal with speaker-specific variability and second, to what extent do listeners add a further component of variability themselves? Generally, it appears that listeners have no issue dealing with speaker variability and can, for example, still derive the intended focus structure produced by different speakers with different prosodic cues \citep{griceIntegratingDiscretenessContinuity2017}. An interesting question is how they accomplish this. \citet{kurumadaThinkingProbabilisticallyStudy2021} considered three ways in which listeners can potentially make sense of speaker variability: Cue integration, scaling/normalization, and learning/storage. In the cue integration model, listeners resolve issues by considering top-down information in the form of expectations alongside phonetic cues. In the scaling/normalization approach, listeners perceptually normalize or scale phonetic cues across speakers in order to remove inter-individual variability in their mental representation. In the learning/storage model, listeners learn speaker-specific prosodic patterns, which they can then exploit. To provide evidence for either of the three models, \citet{kurumadaThinkingProbabilisticallyStudy2021} reanalyzed the production data first presented in \citet{muckeEffectFocusMarking2014} and \citet{griceIntegratingDiscretenessContinuity2017} with German speakers producing lexical items in narrow and contrastive focus and the perception data on the same material from \citet{cangemiListenerSpecificPerceptionSpeakerSpecific2015}. Results revealed support for the assumption that listeners store speaker-specific distributions of cues to make sense of variability. Listener behavior was thus best represented by the learning/storage model.

Further evidence for such a model comes from different (West-Germanic) languages and prosodic domains. \citet{xieEncodingDecodingMeaning2021} found that American English listeners can benefit from learning speaker-specific cues to questions and statements. The extent of the benefit naturally depended on the degree to which a speaker varied from the group average. Listeners may also be explicitly trained to pay selective attention to one cue over another. \citet{francisCuespecificEffectsCategorization2008} successfully taught American English listeners to use either VOT or onset F0 to identify voicing contrasts in stop consonants. Similarly, Dutch listeners were able to learn speaker-specific cue weightings for lexical stress in \citet{severijnenTrackingTalkerspecificCues2023}. They were exposed to one speaker that used F0 as a cue and another that used intensity, and were able to store and apply this information to ambiguous stimuli. Listeners are thus flexible in the type of cue they predominantly pay attention to. This indicates that perceptual preferences are amenable to experience (see Section \ref{sec: experience} for more details on the effects of phonetic and musical experience on individual variability).

In general, listeners appear to be extremely flexible in prosodic processing. \citet{ipSearchSalienceFocus2021} found that Australian English listeners easily dealt with speaker variability in their use of cues to focus in a preceding word. In a phoneme detection task, they investigated to what extent listeners were able to use prosodic cues in a word preceding the target word that naturally varied across speakers. Listeners were in fact sensitive to preceding prosody irrespective of the number of cues used by the speaker producing the stimulus. The effect did not differ depending on the speaker-specific strength of encoding, no single cue or cue combination was preferred by the listeners.

In addition, listeners are able to generalize across speakers and to some extent disregard speaker-specific variability when categorizing speech sounds \citep{kriengwatanaGeneralAuditoryBias2015, roettgerEvidentialStrengthIntonational2019}. Especially when speakers produce unreliable intonational cues, listeners may learn to ignore them but still make use of such cues in other speakers \citep{roettgerListenersAdaptationUnreliable2020}. In a perception experiment, German listeners were exposed to two different speakers, one of whom marked a given-new contrast intonationally as expected and the other produced a reversed and thus unreliable mapping \citep{roettgerListenersAdaptationUnreliable2020}. After listening to these stimuli, listeners had to identify the correct referent in a forced-choice task. When encountering reliable speakers, listeners were able to use intonation as a cue to select the correct referent. However, when they encountered an unreliable speaker, listeners had learned to disregard the intonational cues produced by that speaker, while still making use of intonation for the other, more reliable speaker. Similarly, \citet{jasminShorttermPerceptualReweighting2022} found that American English listeners quickly learned to discard unreliable durational cues to nuclear accent placement, instead relying solely on F0 modulations that were presented consistently as correlates of accenting. These findings indicate that listeners can adapt to speaker-specific variability, even when encountering unexpected form-meaning mappings.

Taken together, these studies provide sufficient evidence that listeners can deal with speaker-specific variability. However, this may still be costly for the listener and depend on individual characteristics, such as their ability to adapt or their sensitivity to prosody. Listeners thus add another source of variability. Generally, differences in the robustness of focus type encoding seem to have direct implications for the interpretation of the utterances on the side of the listener: Speakers who mark contrasts more clearly are more easily interpretable by the listener, while speakers producing weaker contrasts are less easily interpretable (\citealt{cangemiListenerSpecificPerceptionSpeakerSpecific2015}, but cf. \citealt{ipSearchSalienceFocus2021}). However, it is not only the speaker's ability to encode meaning in the signal that determines communicative success. Listeners also differ in how competently they interpret prosody.
For example, \citet{cangemiListenerSpecificPerceptionSpeakerSpecific2015} found different levels of proficiency in German listeners' ability to distinguish between focus types, which are encoded via prosodic prominence. In line with this finding, \citet{jabeenVariabilityPerceptionProsody2021} observed that listeners differed in their willingness to accept prosody-context mismatches, specifically, between pitch accent type and focus type. While most listeners behaved as expected, in that they preferred H* or L* accents on constituents in broad focus and L+H* accents on constituents in contrastive focus, some listeners showed no preference or even the reverse pattern. The willingness to accept such mismatches may indicate greater flexibility and thus constitute an advantage in dealing with speaker variability as discussed in the preceding paragraph. However, in the previously discussed studies by \citet{roettgerListenersAdaptationUnreliable2020} and \citet{jasminShorttermPerceptualReweighting2022}, listeners specifically learned to ignore productions by unreliable speakers or unreliable cues, while still making use of reliable cues. In the study by \citet{jabeenVariabilityPerceptionProsody2021}, some listeners were generally unreceptive to prosody-context mismatches, which indicates a lack of competence in interpreting prosodic cues to meaning contrasts rather than adaptability.

Perception studies using Rapid Prosody Transcription \citep[RPT,][]{coleNewMethodsProsodic2016} usually find more uniform behavior in the marking of the absence of prominence than the presence of prominence \citep[e.g.,][]{coleNewMethodsProsodic2016}. \citet[181]{laddProsodicProminenceLanguages2023} argue that this discrepancy is related to the fact that ``prominence need not involve a consistent and unified set of acoustic exponents." Individual raters are thus quicker to agree on what does not constitute prominence than what does, which indicates that listeners may naturally gravitate towards different cues that they predominantly attend to when decoding prominence. In this sense, redundancy in the prosodic signal facilitates across-listener differences. 

\newpage
Accordingly, perception studies have observed substantial variability in the cues listeners attend to in decoding prosodic prominence. For example, \citet{royIndividualDifferencesPatterns2017} investigated various prosodic (F0, intensity, duration) and non-prosodic (word frequency, part-of-speech) cues to prominence in an RPT task with American English listeners. They found that individuals differed in how many of these cues they used with a clear hierarchy emerging: Part-of-speech and word frequency were used by the most participants followed by phrase boundary placement, intensity, F0 peak, and post-pause duration. In addition, \citet{royIndividualDifferencesPatterns2017} observed an implicational hierarchy: If listeners used intensity, they also relied on F0, if they used intensity and F0, they also attended to post-pause duration and if they used these three cues, they also made use of word frequency. Interestingly, 15 out of 33 listeners made use of all cues investigated in this study, while there were two listeners for whom prominence judgments were not related to any of the cues. Another worthwhile observation from this study is that there was more variability in the cues listeners attended to when decoding prominence than when decoding phrase boundaries. This indicates that prominence is special in facilitating more listener-specific variability.

In a study on German, \citet{baumannWhatMakesWord2018} identified two groups of listeners: The larger group (about two thirds of the sample) attended primarily to cues related to the F0 contour, while the smaller group seemed to rely mostly on non-prosodic cues such as word frequency, part-of-speech, and argument structure in addition to duration. This finding was confirmed in a similar study on prominence perception in German by \citet{wagnerExploitingSpeechgestureLink2019}, who identified two groups of listeners, one relying mostly on signal-related cues and the other integrating both structural and signal-based cues. In a study on resolving syntactic ambiguities in German, \citet{vanderburghtPitchAccentsCreate2021} similarly observed that some listeners were more sensitive to the prosodic signal, while others relied on their structural biases.

Comparable findings have been reported for individual reliance on prosody and lexical information in two studies on American English. \citet{ishidaPeopleAreMore2016} using a degraded speech signal showed that some listeners relied more on the speech signal and others on prior knowledge/lexical influences in order to restore missing or manipulated phonemes. This was confirmed by \citet{giovannoneIndividualDifferencesLexical2023} in a similar study, where they also observed some evidence for a trade-off effect between acoustics and lexical information. The latter additionally tested participants at two points in time approximately two weeks apart and found that listener performance was similar in both sessions, indicating that the reliance on lexical information is a stable listener trait rather than a variable state.

However, it is possible that such trade-off effects do not (only) arise from listener preferences for certain types of cues, but are related to different task interpretations. \citet{coleListeningSoundListening2014} pursued the idea that American English listeners' attention could be directed to different aspects of prominence, namely acoustic cues versus meaning-based cues, by means of the instructions given in the experiment. They found confirmation for their hypothesis in that listeners differed in their prominence judgments as a function of the instructions they received. Interestingly, given non-specific instructions, prominence judgments resembled those gleaned from the condition in which participants were asked to focus on acoustic cues. In a follow-up study, it was confirmed that the difference in behavior did indeed stem from listeners attending more to phonetic cues (i.e., duration, intensity, and F0) when instructed to listen for prominence and to lexical-pragmatic cues (i.e., word frequency) when instructed to base their decision on meaning \citep{coleSoundStructureMeaning2019}. This effect was larger in French and Spanish listeners than in English listeners, indicating that correspondence between lexical-pragmatic factors and prosody is generally closer in English, according to the authors.

In addition to differences stemming from differing interpretations of experimental tasks, prior expectations may also variably affect listeners. In fact, \citet{moProsodyProductionPerception2010} observed more variability in prominence ratings across listeners in a mock RPT task without any auditory input being presented than in a regular RPT task with audio. This observation indicates that expectations with regard to prominence are more variable than the actual perception of prosodic prominence. Since it has been argued that expectations are crucial in the interpretation of information structure encoded via metrical structure \citep{calhounCentralityMetricalStructure2010}, this is yet another factor adding to inter-individual variability in the perception of prosody.

\citet{roettgerPositionalBiasesPredictive2021} found that American English listeners differed in whether or not they were able to make use of the presence of a prenuclear pitch accent to anticipate the discourse status of an upcoming referent: Some listeners learned to integrate this information, others ignored it. Such findings could contribute to resolving theoretical discussions around topics such as the contribution of prenuclear accents to meaning \citep{buringSemanticsIntonationInformation2007, baumannInfluenceInformativenessProsody2021, roessigInverseRelationPrenuclear2024}.

Considering another modality that has been shown to interact with prominence, \citet{bujokAudiovisualPerceptionLexical2024} tested the influence of audiovisual cues on the perception of lexical stress in Dutch. They note that there was quite some inter-individual variability in the extent to which listeners were able to make use of visual cues in the face of the speaker to determine stress patterns. The addition of a new layer of redundancy thus introduces more opportunity for individual variability.

Individual differences in perception may be linked to production. However, evidence for this relation is not clear-cut. On the segmental level, \citet{beddorTimeCourseIndividuals2018} found evidence for a link in the production and perception of coarticulatory nasalization in American English speakers: Participants who produced earlier nasalization were also more efficient at attending to this cue throughout the signal in perception. In contrast, \citet{schertzIndividualDifferencesPhonetic2015} did not observe a direct relation between the production and perception of different cues (VOT, F0 at vowel onset, and closure duration) to distinguish stop consonants in speakers of Korean, although individual listeners were quite variable in their cue weightings. On the suprasegmental level, \citet{kimIndividualDifferencesProduction2020}, who explicitly investigated individual variability in phrase boundary marking, also failed to provide evidence for a production-perception link. However, \citet{speerImplicitProsodyCorrective2015} observed that individual differences in prosody production predicted sensitivity to priming in a silent-reading task. Listeners were presented with written text that contained a contrast, e.g., ``Jacquelyn didn’t pass the test. Belinda passed the test." followed by an auditory stimulus with matching or mismatching intonation, i.e., ``Belinda" with an L+H* accent (match) or deaccented (mismatch). Participants were asked to indicate via a button-press whether the word they heard was the first word of the second sentence they had read. Results suggested that there was a facilitative effect of the matching condition only in those listeners that produced the same prosodic contrast in a subsequent reading task (i.e., L+H* accents on the corrective item), but not for participants who did not produce such a contrast. 

Beyond a potential production-perception link, there may be personality traits or experiences that shape the way individuals produce and perceive prosody. The next section will explore some of the causes for individual variability that have been suggested.

\section{Causes for individual variability} \label{sec:causes_variability}

There are a number of potential causes underlying the degree and kind of variability that has been observed in prosodic research. These relate to communicative skill, empathy, the sociodemographic identity of the speaker/listener, and experience, especially with prosody or music as a related domain. Research has also been concerned with the role of neurodiversity, e.g., autism spectrum disorder \citep{wehrleConversationIntonationAutism2023}, or developmental disorders \citep{schwartzSpeechPerceptionLexical2013} in speech perception and production. However, the present book will not discuss neurodiversity or clinical conditions as underlying causes for variability, as none of the experiments explicitly target these populations or employ appropriate screening tools.

\subsection{Communicative skill and empathy}

Under listener-oriented accounts of speech production, speakers take into consideration the listener's knowledge and demands in communicative situations and adjust their speech accordingly \citep[e.g.,][]{lindblomExplainingPhoneticVariation1990}. Doing this successfully requires an adequate theory of mind \citep{turnbullListenerorientedPhoneticReduction2019}, which is the ability to impute mental states to oneself and to others \citep{premackDoesChimpanzeeHave1978}. Higher theory of mind abilities should thus facilitate the appropriate use of prosodic prominence to manage listeners' expectations and guide their attention in conversation \citep{turnbullListenerorientedPhoneticReduction2019}. From a perceptual side, individual variability has been explained in terms of different cognitive processing styles \citep{bishopPhonologyPhoneticsSignalextrinsic2020, yuIndividualDifferencesLanguage2019}.
 Related to both theory of mind and cognitive processing styles are the concepts of communicative skill and empathy. 
 
Communicative skill has been measured in different ways, both explicitly and implicitly. A common direct way of measuring communicative skill is via a self-administered questionnaire, namely the Autism-Spectrum Quotient (AQ) questionnaire \citep{baron-cohenAutismSpectrumQuotientAQ2006}. As the name suggests, the AQ questionnaire is intended to quantify autistic traits in adults. Such autistic traits can include inhibited communicative skill. In fact, ten out of the 50 questions in the AQ directly query skills and preferences in conversation. The AQ score is inversely correlated with communicative skill: A higher score indicates higher levels of autistic traits and thus lower communicative skill. Studies in prosody aiming to relate individual variability to communicative skill have used either the full AQ questionnaire \citep[e.g.,][]{turnbullListenerorientedPhoneticReduction2019} or only the ten-question subset related to communication \citep[e.g.,][]{bishopFocusProjectionPrenuclear2017, bishopPhonologyPhoneticsSignalextrinsic2020} to quantify communicative (or sometimes called pragmatic) skill.

More indirect measures, which are more tedious to collect and score, and thus less widespread in use, are the Reading the Mind in the Eyes (RMITE) test and the Strange Stories task. The RMITE test requires participants to infer what people are feeling from a photograph of only their eye region \citep{baron-cohenAnotherAdvancedTest1997}. In the Strange Stories task, participants read a short story and subsequently answer an open question that is supposed to reveal the participant's understanding of the characters' relationships and motivations \citep{happeAdvancedTestTheory1994}. \citet{turnbullListenerorientedPhoneticReduction2019} used both the RMITE test and the Strange Stories task in addition to the AQ questionnaire and found them all to be weakly correlated, indicating that they measure similar abilities.

The AQ has been found useful to explain differences in processing pragmatic information. \citet{nieuwlandIncrementalityPragmaticProcessing2010}, for example, investigated the relation between communicative skill and the processing of scalar implicatures in a silent reading task measuring Event-Related Potentials. They found that their American English-speaking participants with higher communicative skill were more likely to exhibit a pragmatic N400 effect, i.e., to make scalar inferences, than participants with lower communicative skill. \citet{zhaoAreScalarImplicatures2015} found further support for this relation in a Mismatch Negativity Study with native speakers of Chinese: Participants with higher communicative skill were more sensitive to scalar implicatures than those with lower communicative skill.

The hypothesis of a linguistic processing advantage for individuals with higher communicative skill has since been extended to prosodic research. In \citet{bishopPhonologyPhoneticsSignalextrinsic2020}, American English listeners with higher communicative skill were found to perform better in a prominence rating task (RPT) than listeners with lower communicative skill. In the same vein, in \citet{orricoInfluenceEmpathyAutisticlike2023}, British English listeners with higher communicative skill more successfully distinguished between H* and L+H* accents. Listeners with lower communicative skill may thus be at a disadvantage when it comes to prosodic processing.

This may also have consequences for the interpretation of prosody-meaning mapping. \citet{bishopIndividualDifferencesTopdown2016} found that American English listeners with higher communicative skill were more sensitive to a change from broad to narrow focus in a prominence rating task than listeners with lower communicative skill. What is more, higher pragmatic skill as indicated by the RMITE test was correlated with higher sensitivity to durational differences. Similarly, in \citet{bishopFocusProjectionPrenuclear2017}, American English listeners with higher communicative skill were more sensitive to the presence or absence of a prenuclear accent as a cue to focus type.

Attempts to relate communicative skill to prosody production are less frequent and have so far yielded mixed results. \citet{turnbullListenerorientedPhoneticReduction2019} found that American English speakers with higher communicative skill enhanced their speaking style more by producing longer words in a clear speech condition, where they were asked to speak to a hearing impaired person or a non-native speaker. In line with this finding, \citet{bishopBriefReportAutistic2022} observed that American English speakers with higher communicative skill generally produced more acoustically distinct vowel categories than speakers with lower communicative skill. However, this was only true for female speakers. In contrast, \citet{turnbullListenerorientedPhoneticReduction2019} observed that speakers with \textit{lower} communicative skill more strongly encoded semantic-pragmatic predictability prosodically, while second mention reduction was not influenced by communicative skill in any way. In fact, the degree of second mention reduction was similar across speakers. However, in a second experiment controlling for general intelligence levels, the relation between communicative skill and encoding of predictability did not replicate. Turnbull interprets these findings as evidence against a listener-oriented account of speech production in which individual differences could be related to differences in theory of mind.

Another psychological construct that potentially relates to linguistic processing aptitude is empathy. Empathy can be quantified via another self-administered questionnaire, the Empathy Quotient (EQ, \citealt{baron-cohenEmpathyQuotientInvestigation2004}). The EQ correlates inversely with the AQ, i.e., individuals with higher levels of autistic traits generally exhibit lower levels of empathy \citep{baron-cohenEmpathyQuotientInvestigation2004}.

Several studies have found empathy levels to play a role in the perception of prosody-meaning mapping \citep[e.g.,][]{esteve-gibertEmpathyInfluencesHow2020, orricoIndividualEmpathyLevels2020, orricoInfluenceEmpathyAutisticlike2023}. \citet{esteve-gibertEmpathyInfluencesHow2020} investigated French listeners' ability to disambiguate referents based on intonation cues in relation to their level of empathy. They found that listeners with higher empathy levels were more sensitive to intonation than listeners with lower empathy levels. The latter were not able to use intonation as a cue to resolve lexical ambiguities. Similarly, in an RPT task with British English listeners, \citet{orricoInfluenceEmpathyAutisticlike2023} observed a correlation between empathy score and pragmatics: Listeners with a higher empathy score were more likely to mark contrastive target words as prominent than listeners with lower empathy scores. This may indicate a higher sensitivity to pragmatic information in more empathetic individuals.

\subsection{Sociodemographics}

Another source of variability is the sociodemographic background of individual speakers and listeners. Although these characteristics are not the primary research objective in the present book, they are introduced here for completeness, since the sample of participants in the studies to be presented is not entirely uniform with regard to these parameters.

One reason for prosodic variability is the region of origin of speakers/listeners. In a large scale study on regional variability in the British Isles, \citet{grabeIntonationalVariationUrban2004} found several distinctive prosodic patterns across speakers from different regions. For example, speakers of Cambridge and Newcastle English used compression while speakers from Belfast and Leeds used truncation to deal with tonal crowding in words with limited voiced segments. She also found differences in preferences for nuclear pitch accent types across the regions under investigation. Similarly, studies on German have identified preferences for certain pitch accent types across different regions: For example, \citet{petersIntonationDeutscherRegionalsprachen2006} identified H*L (described as a falling accent, corresponding to H* in the GToBI system) as the most common nuclear pitch accent in Northern German varieties and L*H (late rise, analogous to GToBI's L*+H) as the most frequent accent in Southern German varieties. Regional differences in peak alignment within the same pitch accent type have also been observed. In prenuclear L*+H accents, Southern German speakers produce later alignment than Northern German speakers \citep{attererPhoneticsPhonologySegmental2004}. These differences could be indicative of different pitch accent inventories. In a study by \citet{zahner-ritterThreeKindsRisingFalling2022}, Northern German speakers more reliably distinguished between three types of pitch accents L+H*, L*+H, and the newly proposed (LH)*, while Southern German speakers seemed to distinguish only L+H* and L*+H accents.

Regional variability can also affect prosody perception, although perception studies investigating this factor are somewhat rare: For example, Glasgow and Belfast English listeners exhibited reduced sensitivity to prominence differences compared to Standard Southern British English listeners \citep{jeonPerceptionPitchHeight2024}. However, \citet{zahner-ritterThreeKindsRisingFalling2022} did not find regional variability in a meaning-attribution task, where the nuclear accent represented pitch contours concomitant with L*+H, (LH)*, and L+H* accents. This result seems to indicate that form-meaning mapping is similar across Northern and Southern German speakers. 

Another oft-cited source of variability is gender. \citet{haanMaleVsFemale1999} observed that female speakers of Dutch produced larger pitch ranges than male speakers. Similarly, \citet{dalyPitchingItDifferently2001} found that female speakers of New Zealand English generally used larger pitch ranges and more dynamic pitch than male speakers in both story telling and sentence list reading tasks. In the data by \citet{clopperEffectsGenderRegional2011}, female speakers of American English used relatively more L*+H accents and fewer H* accents than male speakers when reading a short passage. In \citet{niebuhrAreThereShapers2011}, summarized in Section \ref{sec:ind_production} above, four out of five ``aligners" were male while three out of five ``shapers" were female, although this is of course a rather weak relation in a very small sample. All of these findings indicate that women rely more on F0 scaling in their productions than men. 

On the perception side, \citet{bishopPhonologyPhoneticsSignalextrinsic2020} observed that male listeners were less proficient at identifying prominence than female listeners in an RPT task. They attribute this to potential differences in communicative skill, which is generally higher in women than men \citep{baron-cohenAutismSpectrumQuotientAQ2001}. In a delayed-recall task, \citet{kochContrastiveIntonationEffects2021} also found that female participants were more sensitive to prosodic focus marking than male participants: Contrastive focus facilitated recall of alternatives only in women.

Generally, women seem to be more sensitive to intonation in perception and more expressive in production (at least in terms of F0 scaling) than men. However, it is possible that these differences are related to personality traits such as pragmatic skill that are unevenly distributed across the genders. In any case, most gender effects seem to be small if detected at all. \citet{peppeProsodicVariationSouthern2000} in a comprehensive study involving 90 British English speakers from London observed no gender effects in prosody production and perception in a series of different tasks. The same study did not find any effects of age, another sociodemographic factor that appears to be studied to a much lesser extent than gender and region of origin.

\subsection{Experience} \label{sec: experience}

Experience plays a role in prominence processing, especially when it comes to phonetic training. Trained experts may behave differently than naive participants, especially in controlled laboratory studies that tap into metalinguistic knowledge \citep[e.g.,][]{laddInfluenceAdjacentPitch1994, royIndividualDifferencesPatterns2017, wagnerExploitingSpeechgestureLink2019}. In \citet{royIndividualDifferencesPatterns2017}, trained annotators exhibited stronger effects on prominence ratings for many prosodic cues in comparison to untrained annotators. Similarly, \citet{wagnerExploitingSpeechgestureLink2019} observed that expert annotators relied more on structural cues to prominence (i.e., stressability of a syllable and part-of-speech) 
than naive listeners when rating prosodic prominence. These effects may be stronger in such meta-linguistic tasks than in more natural settings and tasks.

Individuals also differ with regard to their sensitivity to pitch changes, which may also be modulated by phonetic training. For example, phonetically trained individuals performed significantly better than untrained individuals in a forced choice task where they had to indicate the larger pitch change out of two \citep{thartDifferentialSensitivityPitch1981}. A related domain that may have similar effects is musical training. \citet{schonMusicSpeechMusic2004} using Event-Related Potentials observed that professional (French) musicians were more sensitive to small pitch changes in music and in language than non-musicians. A similar finding has been reported by \citet{thartDifferentialSensitivityPitch1981}.

However, some research indicates that a distinction is to be made between musical training and musical aptitude. For example, \citet{tierneyIndividualDifferencesPerception2021} found that musical aptitude but not musical training predicted the strength of a speech-to-song illusion in listeners, the speech-to-song illusion being the perceptual phenomenon that a spoken phrase presented multiple times is eventually heard as sung. According to this finding, musicality is rather an individual trait and not so much dependent on experience.

\citet{niebuhrF0basedRhythmEffects2009} argues that self-proclaimed more musical individuals may have an advantage at doing more analytic, metalinguistic tasks, which however does not indicate a linguistic benefit. Evidence for this claim comes from the comparison of the results of two different studies, both testing the effects of rhythmic context on prominence detection. In \citet{niebuhrEffectGlobalRhythms2008}, participants were asked to identify the word that perceptually stands out and an effect of the rhythmic manipulation was found only for the group of musical participants. In contrast, \citet{niebuhrF0basedRhythmEffects2009}, using a reproduction task in which participants had to repeat what they previously heard, found no difference in the observed effect of rhythmic context on prominence perception between the two groups. These findings suggest that some differences may be due to the nature of the task, which can benefit individuals with certain skills more than others.

More generally, \citet{peppeProsodicVariationSouthern2000} found a correlation between education level and sensitivity to prosody in different perception tasks. Participants receiving education beyond the age of 18 scored higher in half of the prosodic performance tasks. This finding may likewise reflect better meta-linguistic awareness or a better understanding of task instructions.

In the experimental part of this book, all three studies will briefly report tendencies regarding sociodemographic information of participants and the studies presented in Chapters \ref{ch:Production II} and \ref{ch: Perception} also take a look at a measure of communicative skill in production and perception, respectively. However, the primary concern of the present work is the systematic description of individual variability in the prosodic encoding and decoding of prominence relations. The following chapter will present a first exploration of individual differences in information status marking.
